% !TEX root = ../tecnologia-en-marcha.tex
% !TEX encoding = UTF-8 Unicode
% Texto: primer borrador del artículo, con acentos restaurados.

\section{Conclusiones y recomendaciones}

El marco operativo desarrollado permitió traducir tres principios de la
Estrategia Nacional de IA de Costa Rica en nueve requerimientos verificables
para sistemas de IA basados en sensores PIR. La implementación demostró que esos
requerimientos pueden evaluarse mediante un procedimiento reproducible de seis
pasos, articulado con el AI Act de la Unión Europea y el NIST AI RMF.

La aplicación sobre dos pilotos del conjunto CASAS generó expedientes
completos, con 18\,736 inferencias registradas en total y cobertura de los
nueve requerimientos en ambos casos. Los dos sistemas evaluados obtuvieron
veredicto de no cumplimiento, lo cual confirma la utilidad del marco para
producir evidencia incluso cuando el resultado es desfavorable para el sistema
auditado.

El estudio aporta una forma concreta de pasar de principios normativos a
evidencia técnica. En particular, integra explicabilidad local y global,
medición de equidad con umbrales declarados previamente, trazabilidad por
inferencia y documentos auditables. Esta combinación permite que un tercero
revise no solo los resultados, sino también el protocolo y los artefactos que
los sostienen.

Como trabajo futuro, se recomienda validar los umbrales de equidad con un panel
multidisciplinario, incorporar pruebas de robustez ante entradas ruidosas o
degradadas, evaluar reproducibilidad de resultado sobre múltiples semillas y
extender la verificación a escenarios de monitoreo continuo. También se
recomienda aplicar el marco a datasets con metadatos demográficos más
completos, siempre que su uso sea éticamente justificable y compatible con la
protección de datos personales.
